% GENERATED FILE. Edit canonical lessons, then run npm run generate:books.
% canonical-source-hash: fnv1a64:ce0efcf0cecd62c7
% canonical-lessons: MR-C45-barobar, MR-C45-pron-oblique, MR-C45-shii, MR-C45-ne, MR-R45-with

\chapter{With Whom, and With What}
\label{ch:mr-with}
\hlchaptermodality{45}

\begin{chapteropening}
\textbf{By the end of this chapter:} \emph{I can say who I am with, who I am speaking to and what I am doing something with, choosing between three endings English collapses into one, and I can put a pronoun in front of any of them.}

Complete MR-R45-with and put all three endings on mitra, saying for each one whether it takes a companion, a partner in the act, or a tool.
\end{chapteropening}

\section[barobar]{\mr{बरोबर} (barobar) — with}

\label{lesson:MR-C45-barobar}



\begin{quote}
You can say where a thing is. You cannot yet say who is there with you — the book has taught \textbf{\mr{मित्र}} since chapter eighteen with no way to attach him to anything.

\begin{itemize}
\raggedright
  \item \emph{Recall:} say \emph{ḍāvā}, then \emph{ujvā}, then \emph{lāmb}
\end{itemize}
\end{quote}

\subsection*{You'll want to know: \mr{बरोबर}}
\begin{quote}
\textbf{\mr{बरोबर}} — \emph{barobar} — \textbf{with, together with}
\end{quote}

\begin{quote}
\textbf{\mr{मी} \mr{मित्राबरोबर} \mr{जातो}.} — \emph{mī mitrābarobar jāto.} — \textbf{I go with my friend.}
\end{quote}

Another postposition, and the same oblique stem you have been bending for four chapters: \textbf{\mr{मित्र} $\to$ \mr{मित्रा}-}, then \textbf{\mr{बरोबर}}. Nothing new about the noun.

The word has a second life you should meet now, because you will hear it far more often than the first. Said alone, \textbf{\mr{बरोबर}} means \textbf{that's right} — agreement, confirmation, \emph{correct}. Look at the base it shares with \textbf{\mr{बरं}}, the \emph{fine, good} of chapter four: \textbf{\mr{बरोबर}} is \emph{good-alongside}, and a thing that lines up alongside another is both \emph{with} it and \emph{right}.

\begin{quote}
\textbf{\mr{बरोबर}!} — \emph{barobar!} — \textbf{that's right!}
\end{quote}

So one word gives you a way to travel with somebody and a way to agree with them.

\subsection*{Your turn}
Say these aloud:

\begin{itemize}
\raggedright
  \item \emph{barobar} — with
  \item \emph{mī mitrābarobar jāto}
  \item \emph{barobar!} — on its own, meaning \emph{that's right}
  \item \emph{mitra tithe āhe} — and then put yourself beside him
\end{itemize}

\begin{tcolorbox}[breakable,colback=teal!4,colframe=teal!35!black,title={Before you move on}]
What does \textbf{\mr{बरोबर}} mean when it stands alone? (\emph{\textbf{That's right}}.) What does \textbf{\mr{मित्र}} become in front of it? (\textbf{\mr{मित्रा}-}.)

Sources: \href{https://en.wiktionary.org/wiki/\%E0\%A4\%AC\%E0\%A4\%B0\%E0\%A5\%8B\%E0\%A4\%AC\%E0\%A4\%B0}{Wiktionary: \mr{बरोबर}}.
\end{tcolorbox}
\section[mājhyābarobar]{\mr{माझ्याबरोबर} (mājhyābarobar) — a pronoun bends too}

\label{lesson:MR-C45-pron-oblique}



\begin{quote}
\emph{With my friend} works. Try \emph{with me} and you will find you cannot say it: \textbf{\mr{मी} \mr{बरोबर}} is not Marathi.
\end{quote}

\begin{grammarlens}[title={Grammar Lens: the pronoun goes through its possessive first}]
A noun takes an oblique stem before a postposition. A pronoun does too, and its oblique stem is built on the \textbf{possessive} you already know rather than on the subject form:

\begin{quote}
\textbf{\mr{मी}} $\to$ \textbf{\mr{माझ्या}-} $\to$ \textbf{\mr{माझ्याबरोबर}} — \emph{with me}
\end{quote}

\begin{quote}
\textbf{\mr{तू}} $\to$ \textbf{\mr{तुझ्या}-} $\to$ \textbf{\mr{तुझ्याबरोबर}} — \emph{with you}
\end{quote}

\begin{quote}
\textbf{\mr{तुम्ही}} $\to$ \textbf{\mr{तुमच्या}-} $\to$ \textbf{\mr{तुमच्याबरोबर}} — \emph{with you} (respectful)
\end{quote}

Read the middle column. Each one is \textbf{\mr{माझं} / \mr{तुझं} / \mr{तुमचं}} with its ending changed to \textbf{-\mr{या}}, which is exactly what a noun in \textbf{-\mr{ए}} did back in chapter forty-one. The machinery is the same; only the starting form is different, and the starting form is the possessive.

That is the rule worth holding, because it saves you from guessing every time: \textbf{to put a pronoun in front of a postposition, take its possessive and bend that.}
\end{grammarlens}

\subsection*{Your turn}
Say these aloud:

\begin{itemize}
\raggedright
  \item \emph{mājhyābarobar} — with me
  \item \emph{tumchyābarobar} — with you, respectfully
  \item which form each one started from — \emph{\textbf{\mr{माझं}}}, \emph{\textbf{\mr{तुमचं}}}
\end{itemize}

\begin{tcolorbox}[breakable,colback=teal!4,colframe=teal!35!black,title={Before you move on}]
Which form of a pronoun does the oblique start from? (\textbf{The possessive}.) So what is \emph{with you}, said to a stranger? (\textbf{\mr{तुमच्याबरोबर}}.)
\end{tcolorbox}
\section[-śī]{-\mr{शी} (-śī) — with, to}

\label{lesson:MR-C45-shii}



\begin{quote}
\textbf{\mr{बरोबर}} puts two people side by side. It does not work when the second person is the one you are TALKING to.
\end{quote}

\subsection*{You'll want to know: -\mr{शी}}
\begin{quote}
\textbf{-\mr{शी}} — \emph{-śī} — \textbf{with, to}
\end{quote}

\begin{quote}
\textbf{\mr{मी} \mr{मित्राशी} \mr{बोलतो}.} — \emph{mī mitrāśī bolto.} — \textbf{I speak to my friend.}
\end{quote}

\textbf{-\mr{शी}} is glued straight onto the oblique stem with no space, which is the first difference from \textbf{\mr{बरोबर}}. The second is the job. \textbf{\mr{बरोबर}} is \emph{alongside} — two people going somewhere together. \textbf{-\mr{शी}} is \emph{engaged with} — the partner a verb reaches toward.

That is why \textbf{\mr{बोलणे}} takes it. English says \emph{speak TO}; Marathi says \textbf{\mr{मित्राशी}}, and the choice is not free: \textbf{\mr{मित्राबरोबर} \mr{बोलतो}} would mean you and your friend were both speaking, side by side, to somebody else.

Keep the pair as a contrast rather than as two words:

\begin{quote}
\textbf{\mr{मित्राबरोबर}} — alongside my friend.
\end{quote}

\begin{quote}
\textbf{\mr{मित्राशी}} — engaged with my friend.
\end{quote}

\subsection*{Your turn}
Say these aloud:

\begin{itemize}
\raggedright
  \item \emph{mitrāśī} — with, to my friend
  \item \emph{mī mitrāśī bolto}
  \item the two side by side — \emph{mitrābarobar}, \emph{mitrāśī}
\end{itemize}

\begin{tcolorbox}[breakable,colback=teal!4,colframe=teal!35!black,title={Before you move on}]
Which of the two does \textbf{\mr{बोलणे}} take, and why? (\textbf{-\mr{शी}} — the person is engaged with, not merely alongside.) Which one is written with no space? (\textbf{-\mr{शी}}.)

Sources: \href{https://en.wiktionary.org/wiki/\%E0\%A4\%B6\%E0\%A5\%80}{Wiktionary: \mr{शी}}.
\end{tcolorbox}
\section[-ne]{-\mr{ने} (-ne) — by, with (an instrument)}

\label{lesson:MR-C45-ne}



\begin{quote}
Two ways of being \emph{with} somebody, and still no way of saying what you did it \textbf{with}.
\end{quote}

\subsection*{You'll want to know: -\mr{ने}}
\begin{quote}
\textbf{-\mr{ने}} — \emph{-ne} — \textbf{by, with (an instrument)}
\end{quote}

\begin{quote}
\textbf{\mr{मी} \mr{डोळ्याने} \mr{पाहतो}.} — \emph{mī ḍoḷyāne pāhto.} — \textbf{I see with my eye.}
\end{quote}

English overloads one word: \emph{with my friend}, \emph{with my eye}. Marathi refuses to. A companion takes \textbf{\mr{बरोबर}} or \textbf{-\mr{शी}}; a \textbf{tool} takes \textbf{-\mr{ने}}, and the two sets never swap. Ask yourself which one the thing is — a person you are with, or a thing you are using — and the ending picks itself.

\textbf{-\mr{ने}} glues to the oblique stem like everything else in this chapter, and a pronoun in front of it bends the same way: \textbf{\mr{माझ्याने}}, \emph{by me}.

One thing to file away rather than learn now. This same \textbf{-\mr{ने}} has a second and much larger job in Marathi: it also marks the SUBJECT of a past transitive verb, which is a rule no European language has. That is a later chapter. For now, note only that the ending you have just met is one you will meet again wearing a different hat.

\subsection*{Your turn}
Say these aloud:

\begin{itemize}
\raggedright
  \item \emph{ḍoḷyāne} — with the eye
  \item \emph{mī ḍoḷyāne pāhto}
  \item \emph{mājhyāne} — by me, with the pronoun bent
  \item the three endings of this chapter and what each takes — \emph{barobar}, \emph{-śī}, \emph{-ne}
\end{itemize}

\begin{tcolorbox}[breakable,colback=teal!4,colframe=teal!35!black,title={Before you move on}]
Which of this chapter's three endings takes a TOOL rather than a person? (\textbf{-\mr{ने}}.) What other job does it have, waiting in a later chapter? (\textbf{Marking the subject of a past transitive verb}.)

Sources: \href{https://en.wiktionary.org/wiki/\%E0\%A4\%A8\%E0\%A5\%87}{Wiktionary: \mr{ने}}.
\end{tcolorbox}
\section[barobar · -śī · -ne]{Ask one question and the ending picks itself}

\label{lesson:MR-R45-with}



\begin{quote}
English says \emph{with} three times and means three things. Name the three Marathi endings that keep them apart.
\end{quote}

\subsection*{Your turn}
\begin{itemize}
\raggedright
  \item \emph{Say it:} \emph{mitrābarobar} — alongside
  \item \emph{Say it:} \emph{mitrāśī} — engaged with
  \item \emph{Say it:} \emph{ḍoḷyāne} — by means of, which no person takes
  \item \emph{Say it:} \emph{mājhyābarobar} — and notice the pronoun bent from its possessive
  \item \emph{Recall:} ask \emph{koṇ tumchyābarobar āhe?} — who is with you?
\end{itemize}

\begin{tcolorbox}[breakable,colback=teal!4,colframe=teal!35!black,title={Before you move on}]
One question sorts all three: \textbf{is this a companion, a partner in the act, or a tool?} Which ending does the tool take? (\textbf{-\mr{ने}}.) Which form of a pronoun feeds all three? (\textbf{The possessive, bent to -\mr{या}}.)
\end{tcolorbox}
